\section{Overview}
\label{sec:overview}

\subsection{Aim}

The aim of the project is to implement the Retrieval-Augmented Generation
(RAG) technique~\cite{lewis2020rag} by building a Python-driven pipeline that
answers questions grounded in uploaded documents, giving back accurate,
source-cited responses.

A RAG system answers in two steps. It first retrieves the passages of a
document collection that bear on a question, then asks a language model to
answer from those passages alone. Grounding the answer in retrieved text ties
it to the documents the user supplied rather than to whatever the model
absorbed during training, and it makes every answer checkable: each claim can
be traced to a passage the reader can open.

The system is aimed at densely cross-referenced documents such as legislation,
regulatory policies and standards. The EU Artificial Intelligence
Act~\cite{euaiact} and the Canadian Tri-Council Policy Statement on research
ethics~\cite{tcps2} served as the main development corpora, alongside other
regulations. Documents of this kind share two properties that a generic RAG
pipeline handles poorly.

\begin{itemize}
  \item \textbf{Questions name locations.} A reader asks about ``paragraph 2
  of Article 7'' or ``chapter 3, page 48'' rather than describing content.
  To a similarity search, ``Article 7'' is two tokens: it retrieves passages
  that \emph{mention} Article 7, which are rarely Article 7 itself.
  \item \textbf{Meaning is distributed.} A provision often states only part of
  a rule and delegates the rest: ``a measure requested under Article 93'',
  ``the condition under paragraph 1, point (b)''. An answer taken from one
  passage in isolation can be literally correct and still misleading.
\end{itemize}

\subsection{Requirements}

Three requirements follow from the aim and from these properties.

\begin{description}
  \item[R1 Grounding.] Answers are drawn only from passages of the uploaded
  documents. When the passages do not contain the answer, the system says so
  instead of answering from general knowledge.
  \item[R2 Citation.] Every passage carries its provenance: the document, the
  page number a reader sees printed on the page, and the passage's position in
  the document's hierarchy, written the way a reader would cite it
  (``Article 101(1), second subparagraph'').
  \item[R3 Structure.] A question that names a part of a document retrieves
  exactly that part, and the provisions it cites can be followed from the
  answer.
\end{description}

Two constraints shaped the implementation. First, the system runs entirely on
local, consumer-grade hardware---the development machine has an NVIDIA
GTX~1650 with 4\,GB of video memory and 16\,GB of RAM---so no document leaves
the machine and no query incurs a fee. Second, nothing may be specific to one
legal system: structural parsing and query handling are written against
general numbering conventions, not against the vocabulary of one document.

\subsection{Architecture}

\Cref{fig:architecture} shows the two paths through the system.
\emph{Ingestion} turns an uploaded file into labelled passages and a graph of
cross-references; \emph{answering} turns a question into retrieved passages and
a generated answer. Both run in a Python backend (\texttt{backend/}) exposed as
an HTTP API, and a desktop application (\texttt{frontend/}) presents the
results. \Cref{tab:packages} maps the backend packages to their
responsibilities and \cref{tab:stack} lists the technologies used.

\begin{figure}[ht]
\centering
\begin{tikzpicture}[
  box/.style={draw, rounded corners=2pt, align=center, font=\footnotesize,
              text width=21mm, minimum height=10mm, inner sep=2pt},
  store/.style={box, fill=gray!15},
  flow/.style={-{Stealth[length=2mm]}, thick},
  feed/.style={-{Stealth[length=2mm]}, dashed},
  lane/.style={rotate=90, font=\footnotesize\bfseries},
]
  \node[lane] at (-1.6,0)    {Ingestion};
  \node[lane] at (-1.6,-2.2) {Index};
  \node[lane] at (-1.6,-4.4) {Answering};

  \node[box]   (upload)   at (0,0)       {Upload\\PDF, DOCX, TXT};
  \node[box]   (clean)    at (3.1,0)     {Read \& clean\\page furniture removed};
  \node[box]   (parse)    at (6.2,0)     {Structure\\parser};
  \node[box]   (chunk)    at (9.3,0)     {Chunker\\labels and pages};
  \node[box]   (embed)    at (12.4,0)    {Embedder\\MiniLM on CPU};

  \node[store] (edges)    at (6.2,-2.2)  {Edge store\\SQLite};
  \node[box]   (refs)     at (9.3,-2.2)  {Reference\\extractor};
  \node[store] (chroma)   at (12.4,-2.2) {Vector store\\ChromaDB};

  \node[box]   (question) at (0,-4.4)    {Question};
  \node[box]   (route)    at (3.1,-4.4)  {Router\\exact or hybrid};
  \node[box]   (retrieve) at (6.2,-4.4)  {Retrieval and\\neighbourhood};
  \node[box]   (select)   at (9.3,-4.4)  {Selection\\context budget};
  \node[box]   (generate) at (12.4,-4.4) {Generator\\Llama 3.2};

  \draw[flow] (upload) -- (clean);
  \draw[flow] (clean) -- (parse);
  \draw[flow] (parse) -- (chunk);
  \draw[flow] (chunk) -- (embed);
  \draw[flow] (embed) -- (chroma);
  \draw[flow] (chroma) -- (refs);
  \draw[flow] (refs) -- (edges);

  \draw[flow] (question) -- (route);
  \draw[flow] (route) -- (retrieve);
  \draw[flow] (retrieve) -- (select);
  \draw[flow] (select) -- (generate);

  \draw[feed] (chroma) -- (retrieve);
  \draw[feed] (edges) -- (retrieve);
\end{tikzpicture}
\caption{The ingestion and answering paths. Shaded nodes are persistent
stores; dashed arrows are reads from them during answering. The reference
extractor reads the stored passages back in reading order, so the edges it
produces use exactly the labels the passages carry.}
\label{fig:architecture}
\end{figure}

\begin{table}[ht]
\centering
\small
\begin{tabularx}{\textwidth}{@{}lX@{}}
\toprule
Package & Responsibility \\
\midrule
\texttt{ingestion}  & Reading PDF, DOCX and TXT files; cleaning; structural
                      parsing; chunking; printed page numbers; reference
                      extraction \\
\texttt{embedding}  & Sentence embeddings and the ChromaDB vector store,
                      including metadata lookups \\
\texttt{retrieval}  & Query parsing, routing to exact lookup, hybrid dense and
                      lexical search, passage selection \\
\texttt{generation} & The generator interface and its Ollama-backed
                      implementation \\
\texttt{graph}      & Cross-reference edges, their SQLite store and
                      neighbourhood queries \\
\texttt{pipeline}   & The two entry points, ingestion and answering; context
                      budget; deterministic replies \\
\texttt{api.py}     & FastAPI endpoints \\
\texttt{telemetry.py} & CPU, GPU and memory readings shown in the interface \\
\bottomrule
\end{tabularx}
\caption{Backend packages.}
\label{tab:packages}
\end{table}

\begin{table}[ht]
\centering
\small
\begin{tabularx}{\textwidth}{@{}llX@{}}
\toprule
Layer & Technology & Reason for the choice \\
\midrule
Pipeline        & Python 3.12              & The reference ecosystem for text
                                             processing, embeddings and model
                                             runtimes. \\
API             & FastAPI, Uvicorn         & Typed request models; the
                                             interface and the tests call the
                                             same endpoints. \\
PDF reading     & PyMuPDF, pypdf           & PyMuPDF exposes layout
                                             geometry---line positions and
                                             font sizes---which footer removal
                                             and page numbering rely on; pypdf
                                             is the fallback. \\
DOCX reading    & python-docx              & Paragraph-level access to Word
                                             documents. \\
Chunking        & LangChain text splitter  & Recursive splitting on paragraph,
                                             line, sentence and word
                                             boundaries. \\
Embeddings      & all-MiniLM-L6-v2         & 384-dimensional sentence
                                             embeddings~\cite{reimers2019sbert,wang2020minilm},
                                             fast enough on a CPU. \\
Vector store    & ChromaDB                 & Embedded and persistent; filters
                                             on metadata inside the store. \\
Lexical scoring & rank-bm25                & Okapi
                                             BM25~\cite{robertson2009bm25}
                                             without an external search
                                             engine. \\
Language model  & Llama 3.2 (3B), Ollama   & Runs locally within 4\,GB of video
                                             memory. \\
Graph store     & SQLite                   & Part of the Python standard
                                             library; no server to run. \\
Desktop shell   & Tauri 2                  & Uses the system web view, so the
                                             application stays small. \\
Interface       & React 19, TypeScript     & Typed components, built with Vite
                                             and styled with Tailwind CSS. \\
Telemetry       & psutil, NVML             & CPU, memory and GPU readings. \\
\bottomrule
\end{tabularx}
\caption{Technology stack.}
\label{tab:stack}
\end{table}

\subsection{Development timeline}

\Cref{tab:timeline} summarises how the system was built. Development
proceeded from the bottom of the pipeline upwards: ingestion first, then
retrieval and generation, then structure and cross-references, with the
interface growing alongside.

\begin{table}[ht]
\centering
\small
\begin{tabularx}{\textwidth}{@{}lX@{}}
\toprule
Period & Milestones \\
\midrule
April--May 2026      & Ingestion layer: readers for PDF, DOCX and TXT behind
                       one loader; fixed-size chunking; unit tests on a PDF
                       fixture. Local embedding layer on ChromaDB. \\
June 2026            & Generator interface and the first end-to-end pipeline,
                       retrieving by embedding similarity alone. Repository
                       split into \texttt{backend/} and \texttt{frontend/}.
                       Endpoints for health, upload and query; chat pane wired
                       to the backend; text cleaning; generation through
                       Ollama. \\
June--July 2026      & Queries scoped to one document; token usage shown per
                       answer; hybrid retrieval with BM25 and reciprocal rank
                       fusion; retrieval returns chunk identifiers and
                       metadata. \\
August 2026          & Structural markers and page boundaries kept during
                       ingestion; queries naming a provision routed to exact
                       metadata lookup; live CPU and GPU telemetry; passage
                       locations shown in the interface. \\
Early September 2026 & Passages ordered against the lost-in-the-middle
                       effect; embedder moved to the CPU; symbol-font glyphs
                       stripped; past answers replayable; per-document
                       statistics; structural units detected generically;
                       answers split into a direct answer and a supporting
                       detail. \\
Mid September 2026   & Cross-reference extraction and edge building; paragraph
                       and point labels; running headers and footers removed;
                       edges stored and a neighbourhood returned with each
                       answer; continuous integration; an interactive
                       cross-reference graph; desktop window sizing and native
                       translucency. \\
Late September 2026  & Graph view replaced by a reference tree; retrieval of
                       the exact subdivision named, with typo tolerance;
                       commands and query cancellation; printed page numbers;
                       passage count decided by the system; text that closes a
                       list recognised as a subparagraph; plain-language
                       summaries for bare references. \\
\bottomrule
\end{tabularx}
\caption{Development timeline.}
\label{tab:timeline}
\end{table}

\subsection{Structure of this report}

Part~I covers the system up to the cross-reference graph: ingestion
(\cref{sec:ingestion}), retrieval (\cref{sec:retrieval}), generation
(\cref{sec:generation}) and the interface (\cref{sec:interface}). Part~II
covers what it took to make the system output cross-references: parsing
structure below the level of a provision (\cref{sec:structure}) and the graph
built on it, including how it is presented (\cref{sec:graph}).
\Cref{sec:status} closes with the current status, measured behaviour and
limitations.
